\chapter{State-space models and Mamba}

\section{The state-space model}

A state-space model reads a sequence and keeps a small hidden state $h$ that summarises everything seen so far
(Figure~\ref{fig:ssm}). In continuous time the state evolves as $h'(t) = A\,h(t) + B\,x(t)$ and the output is read
from it as $y(t) = C\,h(t) + D\,x(t)$. To apply this to tokens, the equation is discretised with a step size
$\Delta$: the state then decays by $\bar A = \exp(\Delta A)$ and takes in $\bar B x$ at every step. S4~\cite{gu2022s4}
showed that with a structured $A$ such models capture very long sequences, and that the same computation can be run
either as a recurrence, one token at a time, or as one long convolution.

\widefig{fig_ref_ssm}{From a continuous state-space model to a selective scan, after Gu et al.~\cite{gu2022s4} and
Gu and Dao~\cite{gu2023mamba}. (a) The continuous system. (b) Discretisation with step $\Delta$. (c) The recurrent
view: one state update per token, at linear cost. (d) Selection: $\Delta$, $B$ and $C$ are computed from each
token, so the model decides what to keep.}{fig:ssm}

Mamba~\cite{gu2023mamba} made the step size and the projections depend on the input. This one change is the whole
idea, and panel (d) of Figure~\ref{fig:ssm} spells out why it matters. When a token produces a small $\Delta$, the
decay $\bar A$ is close to one and the input term close to zero, so the state is held and the token is ignored.
When it produces a large $\Delta$, the past decays away and the state is rewritten from the token. The model can
therefore learn, token by token, what to remember and what to skip.

\section{The Mamba block and its two-dimensional form}

Figure~\ref{fig:mamba} draws the Mamba block beside its adaptation to images. The block expands the channels,
mixes neighbouring tokens with a short convolution, runs the selective scan, and multiplies the result by a gate
computed from the input, a structure close to the gated feed-forward layers of Transformers. Blocks are stacked
with a normalisation in front and a residual connection around each.

An image is not a sequence, and flattening it row by row breaks its neighbourhoods. VMamba~\cite{liu2024vmamba}
answers this with a two-dimensional scan, SS2D: the same feature map is read in four orders, along rows and
columns in both directions, each order is scanned separately, and the four results are added back on the grid.
Every pixel then receives context from all four directions at linear cost. Vision Mamba~\cite{zhu2024vim} follows a
related bidirectional design, and MambaVision~\cite{hatamizadeh2025mambavision} shows that pairing the scan with a
plain convolutional branch helps it to keep local detail.

\widefig{fig_ref_mamba}{(a) The Mamba block, after Gu and Dao~\cite{gu2023mamba}. (b) The four-way cross-scan of
VMamba~\cite{liu2024vmamba}, with an independent selective scan (S6) per direction and a cross-merge. (c) The VSS
block that wraps SS2D, after Liu et al.~\cite{liu2024vmamba}.}{fig:mamba}

\section{What \method{} takes from these models}

\method{} borrows four specific pieces and adds one of its own. From Mamba it takes the input-dependent step size.
From VMamba it takes the four-direction scan over a two-dimensional map. From MambaVision it takes a local
convolutional branch beside the scan. From the practice of zero-initialised residual branches it takes an output
projection that starts at zero, so a new block begins as the identity of its residual path. The addition is the
reliability gate: \method{} multiplies each token's step size by how reliable its camera is at that location, which
lets the scan skip the tokens of a camera that cannot be trusted. Chapter~\ref{ch:cffm} draws how.
